\section{Background}\label{sec:background}
\subsection{Ordered set partitions}\label{sec:set_partitions}
We write $[n] = \{1, 2, \dots, n\}$. An \newword{ordered set partition} of $n$ is a sequence $\pi = (\pi_1, \pi_2, \dots, \pi_d)$ of sets where 
\begin{itemize}
    \item each $\pi_i \neq \emptyset$,
    \item $\bigcup_{i} \pi_i = [n]$, and
    \item $\pi_i \cap \pi_j = \emptyset$ if $i \neq j$.
\end{itemize}
We write such an ordered set partition as 
$(\pi_1 \mid \pi_2 \mid \dots \mid \pi_d)$ for visual distinctiveness. The sets $\pi_i$ are called the \newword{blocks} of the ordered set partition. If we forget the ordering of the blocks in an ordered set partition, the result is an \newword{(unordered) set partition} 
$\{\pi_1, \dots, \pi_d \}$. 


We draw a set partition $\pi = \{\pi_1, \dots, \pi_d \}$ of $n$ by placing dots labeled $1, 2, \dots n$ clockwise around the boundary of a disk and then, for each $\pi_i$, drawing the convex hull of the boundary dots whose labels are in $\pi_i$. If these convex hulls do no intersect, we call the set partition a \newword{noncrossing set partition}. We use $\nc(n,d,r)$ to denote the set of a noncrossing set partitions of $n$ with $d$ blocks and blocks size at least $r$.

Let $\ordsetpart(n,d,r)$ denote the set of ordered set partitions of $n$ with $d$ blocks and blocks of size at least $r$. For example, $(2~5~6\mid 3 \mid 1~4)$ and $(3\mid 2~5~6 \mid 1~4)$ are distinct ordered set partitions in $\ordsetpart(6,3,1)$, but they are not distinct when viewed as unordered set partitions. Note that if $r \leq r'$, then $\ordsetpart(n,d,r') \subseteq \ordsetpart(n,d,r)$. We write $\ncop(n,d,r)$ for the noncrossing subset of $\ordsetpart(n,d,r)$.

\begin{figure}[ht]
\includegraphics[scale=.9]{figures/partitionexample.pdf}
\caption{A visual depiction of the ordered set partition $(2~5~6\mid 3 \mid 1~4)$ considered in Section~\ref{sec:set_partitions}. In this picture, the ordering of the blocks is not recorded. In later examples, we will occasionally use color-coding to illustrate the intended ordering, as needed.}
\label{fig:setpartition}
\end{figure}

\subsection{Grassmannians and exterior algebras}\label{sec:Grassmann}
The \newword{Grassmannian} $\Gr(k,n)$ is the parameter space of $k$-dimensional complex vector subspaces of $\C^n$. (In this paper, we will mostly be interested in the case of $\Gr(n,2n)$.) We realize $\Gr(k,n)$ as a smooth projective variety as follows. Consider a $k \times n$ matrix of distinct indeterminates 
\[
M = \begin{bmatrix}
    x_{11} & x_{12} & \dots & x_{1n}  \\
    x_{21} & x_{22} & \dots & x_{2n} \\
    \vdots & \vdots & \ddots & \vdots \\
    x_{k1} & x_{k2} & \dots & x_{kn}
\end{bmatrix}
\]
and let $\C[M]$ be the polynomial ring in these $kn$ variables. Let $R$ be the subring generated by all of the $k \times k$ minors of $M$. Then $R$ is the homogeneous coordinate ring of $\Gr(k,n)$, meaning that we can take $\Gr(k,n) \coloneqq \Proj R$. 

The action of $\SL_k(\C)$ on $M$ by left multiplication induces an action of $\C[M]$. By the Fundamental Theorem of Invariant Theory, the subring $R$ consists exactly of those polynomials in $\C[M]$ that are invariant under this $\SL_k(\C)$ action.

It is often convenient to replace the ring~$R$ by an isomorphic ring~$R'$. Here, $R' = \C[x_{\Delta} : \Delta \in \binom{[n]}{k}] / {\sim}$, where $\binom{[n]}{k}$ denotes the collection of $k$-element subsets of $\{1, 2, \dots, n\}$, each $x_\Delta$ stands in for the corresponding degree $k$ generator of $R$, and ${\sim}$ is the ideal corresponding to the relations among the generating minors of $R$. In this language, the indeterminates $x_\Delta$ are called \newword{Pl\"ucker variables} and elements of ${\sim}$ are called \newword{Pl\"ucker relations}. We will treat these two perspectives interchangeably. 

\subsection{Specht modules}
\label{sec:Spechtbackground}

Consider a $\Z^n$-grading on $\C[M]$, where each variable $x_{ij}$ has degree $\mathbf{e}_j$, where $\mathbf{e}_j$ is the $j$th standard basis vector of $\Z^n$. We are especially interested in the vector subspace $S$ of $R$ spanned by polynomials of multihomogeneous degree $(1,1, \dots, 1) = \sum_{j = 1}^n \mathbf{e}_j$. Note that $S$ consists exactly of multilinear $\SL_k$-invariant functions of the columns of $M$. We will assume that $n = dk$ for some integer $d$, so that $S$ is nontrivial.

It is well-known that $S$ is a finite-dimensional complex vector space whose dimension is given by the \emph{hook-length formula} for \emph{standard Young tableaux} of shape $(d^k)$. Precisely, 
\[
\dim S = \prod_{i=1}^{k} \frac{(d+k-i)!}{(k-i)!}.
\]

The symmetric group $\mathfrak{S}_n$ acts on $M$ by permuting columns, and hence on $\C[M]$. Note that $S$ is an invariant subspace for this action. In fact, it is irreducible as an $\mathfrak{S}_n$-module and is the \newword{Specht module} $S^{(d^k)}$.

In general, there is a Specht module $S^\lambda$ for every integer partition $\lambda = (\lambda_1 \geq \lambda_2 \geq \dots 0)$ with $\sum_i \lambda_i = n$; moreover, every irreducible complex representation of $\mathfrak{S}_n$ is isomorphic to exactly one of these Specht modules. Here, we give an explicit but terse construction of general Specht modules along the lines of the construction of $S^{(d^k)}$ above; for more details of this approach, see \cite[\S2.3]{Patrias.Pechenik.Striker}, while for more standard (and more thorough) textbook treatments, see \eg \cite{Fulton.Harris,Fulton:YoungTableaux,Sagan}. 

Let $\lambda$ be any partition of $n$ and let $\mu$ be the partition whose Young diagram is the transpose of that of $\lambda$. For example, if $\lambda = (d^r,1^{n - r \cdot d})$, then $\mu  = (n - (r-1) \cdot d, r^{d-1})$. Suppose that 
\[
\mu = (\mu_1 \geq \mu_2 \geq \dots \geq \mu_\ell > 0).
\]
Let $M$ be the matrix of indeterminates from Section~\ref{sec:Grassmann} with $k = \mu_1$. Consider the polynomial ring $\mathbb{C}[M]$ in those $n \cdot \mu_1$ variables. For each $\mu_i$, let $\mathsf{S}_i$ denote the set of sequences $\mathbf{j} = (j_1, j_2, \dots, j_{\mu_i})$ with $1 \leq j_1 < j_2 < \dots < j_{\mu_i}$. (Note that if $\mu_i = \mu_{i'}$, then $\mathsf{S}_i = \mathsf{S}_{i'}$.) For $\mathbf{j} \in \mathsf{S}_i$, let $p_{\mathbf{j}} \in \mathbb{C}[M]$ denote the polynomial that is the $\mu_i \times \mu_i$ minor of $M$ involving the top $\mu_i$ rows and the columns indexed by the sequence $\mathbf{j}$. Now, consider the set $\mathsf{S}$ of tuples of sequences $(\mathbf{j}_1, \dots, \mathbf{j}_\ell)$ such that each $\mathbf{j}_i \in \mathsf{S}_i$ and the sets $\mathbf{j}_i$ partition $[n]$ (that is, they are pairwise disjoint with union $[n]$).
For each $(\mathbf{j}_1, \dots, \mathbf{j}_\ell) \in \mathsf{S}$, let \[
p_{(\mathbf{j}_1, \dots, \mathbf{j}_\ell)} \coloneqq \prod_{i=1}^\ell p_{\mathbf{j}_i}.
\]
The vector span of the polynomials $p_{(\mathbf{j}_1, \dots, \mathbf{j}_\ell)}$ for $(\mathbf{j}_1, \dots, \mathbf{j}_\ell) \in \mathsf{S}$ is the Specht module $S^\lambda$. Here, the symmetric group $\mathfrak{S}_n$ is acting by permuting the columns of $M$.

The construction above naturally identifies $S^\lambda$ with a slice of the multihomogeneous coordinate ring of a partial flag variety, extending the construction above of $S^{(d^k)}$ inside the homogeneous coordinate ring of a Grassmannian; for more details, see \cite[\S2.3]{Patrias.Pechenik.Striker}.
One of our main results in Section~\ref{sec:GrassmannCayley} will be a new, and somewhat mysterious, construction of the Specht module $S^{(d^r,1^{n-r\cdot d})}$ inside the homogeneous coordinate ring of the Grassmannian $\Gr(n,2n)$.

\section{Jellyfish tableaux and  invariants}
\label{sec:invariants}
In this section, we define a polynomial associated to each ordered set partition in $\ordsetpart(n,d,r)$; these polynomials will be our web invariants. Although an ordered set partition $\pi \in \ordsetpart(n,d,r)$ is also an element of $\ordsetpart(n,d,r')$ for any $r' \leq r$, the associated web invariant depends on the specified $r$. When $r$ is even, the invariants we define will be independent of the ordering of the blocks; that is to say, they depend only on the underlying unordered set partition. When $r$ is odd, the ordering of the blocks changes the invariant by a predictable global sign, which we discuss in Lemma~\ref{lem:row_col_swap}. For the case $r=2$, our invariants recover those introduced in \cite{Patrias.Pechenik.Striker}.

We begin by defining a set of tableaux that we will use to construct polynomials associated to ordered set partitions. Because the shape consists of $r$ full rows of boxes followed by one box in each additional row and thus resembling jellyfish, we call these $r$-\emph{jellyfish tableaux} (cf. Example~\ref{ex:RStableaux}).

\begin{figure}[ht]
\raisebox{1.4in}{\begin{ytableau}
1 & 3 & 5 & 2\\
4 & 8 & 6 & 7\\
12 & 10 & 9 & 11\\
\none & 16 & \none \\
\none & \none & 13 & \none  \\
15 & \none & \none\\
\none & \none & \none & 14
\end{ytableau}}\hspace{.65in}
%\includegraphics[height=5cm]{figures/cute_jellyfish.jpg}
